\title{xLSTM: Extended Long Short-Term Memory}

\begin{abstract}
The foundational concepts of the Long Short-Term Memory (LSTM) architecture—the constant error carousel and gating mechanisms—were established during the 1990s. Since their inception, LSTMs have demonstrated persistent relevance, playing a crucial role in various deep learning achievements and serving as the basis for the earliest Large Language Models (LLMs). Nevertheless, the introduction of the Transformer, characterized by its highly parallelizable self-attention mechanism, signaled a paradigm shift that enabled superior scalability compared to LSTMs. This work investigates a fundamental question: What is achievable in language modeling if we scale LSTMs to the billion-parameter regime, apply state-of-the-art advancements from modern LLMs, and address the classical deficiencies of LSTMs? To this end, we first propose exponential gating augmented with normalization and stabilization procedures. Next, we revise the LSTM memory configuration and propose: (i) sLSTM, featuring scalar memory, scalar updates, and an updated mechanism for mixing memory; and (ii) mLSTM, which replaces the traditional memory with a matrix representation and employs a covariance-based update, facilitating full parallelization. By incorporating these LSTM enhancements within residual block designs, we construct xLSTM blocks, which can be residually layered to form complete xLSTM systems. The introduction of exponential gating and novel memory mechanisms enables xLSTM to deliver competitive or superior outcomes relative to leading Transformers and State Space Models, both in terms of modeling effectiveness and scalability.\\
Code available at: \url{https://github.com/NX-AI/xlstm}
\end{abstract}

\section{Introduction}
\label{sec:intro}

The Long Short-Term Memory (LSTM) ideas~\citep{Hochreiter:91,Hochreiter:97nips,Hochreiter:97}, i.e., the constant error carousel and gating, were introduced to overcome the vanishing gradient problem of recurrent neural networks~\citep{Hochreiter:91,Hochreiter:00book}:

\begin{align}
\label{eq:lstm_idea}
  c_t \ = \ f_t \ c_{t-1} \ + \ 
          i_t \  z_t \ , \quad  
          h_t \ = \ o_t \ \psi({c_t}) \ . 
\end{align}

The constant error carousel functions by incrementally updating the cell state $c_{t-1}$ (indicated in green) through the addition of cell inputs $z_t$, with these additions being regulated by sigmoid-based gates (depicted in blue). Specifically, the input gate $i_t$ and the forget gate $f_t$ govern this updating mechanism, whereas the output gate $o_t$ manages the emission from the memory cell, namely the hidden state $h_t$. Before producing the hidden state, the cell state undergoes normalization or squashing via $\psi$, after which the output gate determines the final hidden state.

LSTMs have found widespread application across a range of fields \citep{Hochreiter:01,Hochreiter:07,Schmidhuber:15}, dominating text generation tasks up until the introduction of Transformers in 2017~\citep{Vaswani:17}. Their utility has been established in multiple sequence-centric challenges, including but not limited to text synthesis~\citep{Graves:13,Karpathy:15blog}, handwriting production~\citep{Graves:13}, sequence-to-sequence machine translation~\citep{Sutskever:14nips}, program evaluation~\citep{Zaremba:14arxiv}, automated image captioning~\citep{Karpathy:15,Hossain:19}, source code creation~\citep{Karpathy:15blog}, as well as hydrological applications such as rainfall-runoff modeling~\citep{Kratzert:18,Kratzert:19} and flood forecasting models~\citep{Nearing:24}. Within reinforcement learning, LSTMs represent the state-of-the-art among sequence processing architectures, exemplified by their role in models such as AlphaStar for StarCraft II~\citep{Vinyals:17short}, OpenAI Five for Dota 2~\citep{OpenAI:19}, and controllers designed for magnetic confinement in nuclear fusion experiments~\citep{Degrave:22short}. A key strength of LSTMs lies in their capacity for abstraction—skillfully capturing and retaining semantic features in their memory units~\citep{Karpathy:15blog}, as highlighted by the identification of specific neuron types, such as those responsible for number and syntax~\citep{Lakretz:19}, linguistic constructs~\citep{Bau:19}, and sentiment analysis~\citep{Radford:17arxiv}. LSTMs remain integral in prominent, real-world implementations~\citep{Degrave:22short,Nearing:24} and continue to demonstrate lasting relevance.

LSTMs, notwithstanding their considerable achievements, are constrained by three significant drawbacks: (i) Lack of revisability in storage decisions. This shortcoming is illustrated using the {\it Nearest Neighbor Search} task (refer also to Appendix~\ref{sec:appNearestNeighborSearch}): When tasked with identifying the sequence element most similar to a provided reference vector, the model must process the sequence in order and report the attached value at the end. As depicted in the left panel of Figure~\ref{fig:lstmProblems}, LSTM exhibits higher mean squared error for this task. It is unable to substitute a previously stored value in the presence of a better match; in contrast, the xLSTM we propose employs exponential gating to effectively address this challenge. (ii) Constraints on storage capacity; specifically, data are forced to reside in scalar-valued cell states. The {\it Rare Token Prediction} task demonstrates this point. The right panel of Figure~\ref{fig:lstmProblems} plots token prediction perplexity on Wikitext-103~\citep{Merity:17}, segmented by token frequency. LSTM consistently underperforms on infrequent tokens, reflecting its bounded memory resources. Our xLSTM mitigates this by incorporating a memory structure based on matrices. (iii) Sequential computation due to memory entanglement—hidden-to-hidden connections tie the state at each time step to its predecessor, precluding parallel processing.

These challenges faced by LSTMs have largely facilitated the dominance of Transformers~\citep{Vaswani:17} in language modeling. A natural question arises: if we can transcend these limitations and scale LSTMs to the capacity of current Large Language Models, what level of language modeling performance becomes attainable?

\section{Extended Long Short-Term Memory}
\label{sec:xLSTM}

In addressing the shortcomings of LSTM, Extended Long Short-Term Memory (xLSTM) proposes two primary changes to the mechanism outlined in Equation~\eqref{eq:lstm_idea}. The first adjustment is the use of exponential gating; the second involves innovative forms of memory structures. Together, these augmentations yield two notable new variants in the LSTM suite: (i) the sLSTM (Section~\ref{sec:sLSTM}), which incorporates a scalar memory, scalar update operation, and supports memory mixing; and (ii) the mLSTM (Section~\ref{sec:mLSTM}), characterized by a matrix-based memory and an update rule predicated on covariance via an outer product, offering full parallelizability. Both variants leverage exponential gating to improve upon conventional LSTM designs. In the case of mLSTM, to facilitate parallelization, memory mixing---i.e., hidden-hidden recurrent interactions---is eliminated. Both sLSTM and mLSTM are capable of extension to multiple memory cells, with sLSTM supporting inter-cell memory mixing. Additionally, sLSTM may utilize multiple heads, where memory mixing is constrained to cells within individual heads but not between the heads. This introduction of multi-headed sLSTM modules, paired with exponential gating, enables a novel form of memory interaction. For mLSTM, the arrangement of multiple heads and cells is functionally equivalent.

Incorporating these LSTM derivatives within residual block modules yields xLSTM blocks (refer to Section~\ref{sec:xLSTMblock}). Stacking these blocks in a residual fashion gives rise to xLSTM architectures (see Section~\ref{sec:xLSTMarchitecure}). The overall structure and its elements are depicted in Figure~\ref{fig:xlstm_sketch}.

\subsection{Review of the Long Short-Term Memory}
\label{sec:orgLSTM}

The foundational concept behind LSTM, as outlined in~\citep{Hochreiter:91,Hochreiter:97nips,Hochreiter:97}, centers on the scalar memory cell, which serves as both the primary mechanism for computation and storage. This design is intended to mitigate issues related to vanishing gradients~\citep{Hochreiter:91,Hochreiter:00book}, accomplished via the constant error carousel mechanism that governs cell state updates. Three gates regulate the memory cell’s dynamics: the input, output, and forget gates. The introduction of the forget gate was later proposed by \citet{Gers:00}. The LSTM cell’s update equations for timestep $t$ are given as follows:

\begin{align}
c_t \ &= \  f_t \ c_{t-1} \ + \ i_t \ z_t &  & &\text{cell state} \\
h_t  \ &= \ o_t \ \tilde{h}_t \ , & \tilde{h}_t \ &= \ \psi \left( c_t\right)
  &\text{hidden state} \\
z_t \ &= \ \varphi \left( \tilde{z}_t \right) \ , 
  &\tilde{z}_t \ &=  \ \mathbf{w}^\top_{z} \ \mathbf{x}_t \ + \
  r_{z}  h_{t-1} \ + \  b_{z} \ \
  &\text{cell input} \\
i_t \ &= \ \sigma \left( \tilde{i}_t  \right) \ , 
  &\tilde{i}_t \ &= \ \mathbf{w}^\top_{i} \ \mathbf{x}_t \ + \
  r_{i}  \ h_{t-1} \ + \  b_{i} \ \
  &\text{input gate} \\
f_t \ &= \ \sigma \left(  \tilde{f}_t \right) \ , 
  &\tilde{f}_t \ &= \ \mathbf{w}^\top_{f} \ \mathbf{x}_t  \ + \
  r_{f}  \ h_{t-1} \ + \  b_{f} \ \
  &\text{forget gate} \\
o_t \ &= \ \sigma \left( \tilde{o}_t \right) \ , 
  &\tilde{o}_t  \ &= \ \mathbf{w}^\top_{o} \ \mathbf{x}_t \ + \
  r_{o}  \ h_{t-1} \ + \  b_{o} \ \
  &\text{output gate} 
\end{align}

The vectors $\mathbf{w}_{z}$, $\mathbf{w}_{i}$, $\mathbf{w}_{f}$, and $\mathbf{w}_{o}$ denote the sets of weights linking the input $\mathbf{x}_t$ to the cell input, input gate, forget gate, and output gate, respectively. The parameters $r_{z}$, $r_{i}$, $r_{f}$, and $r_{o}$ serve as the recurrent weights mapping from the previous hidden state $h_{t-1}$ to the cell input and the associated gates. Corresponding bias vectors are labeled as $b_{z}$, $b_{i}$, $b_{f}$, and $b_{o}$. The cell input and hidden state activation functions, $\varphi$ and $\psi$, are typically realized as $\tanh$. Specifically, $\psi$ is responsible for bounding the cell state, preventing it from diverging. The activation for all gates employs the sigmoid nonlinearity, given by $\sigma \left( x \right) = 1/(1+\exp(-x))$. Later developments introduced the aggregation of several scalar memory cells $\mathbf{c}_t \in \mathbb{R}^{d}$ into a vector, which in turn permitted the integration of recurrent matrices $\mathbf{R} \in \mathbb{R}^{d \times d}$ to facilitate interaction among different memory cells’ outputs~\citep{Greff:15}. Further explanations are available in Appendix~\ref{sec:apporgLSTM}. Empirical investigations via ablation confirmed that each part of the memory cell architecture is indispensable~\citep{Greff:15}.

\subsection{sLSTM}
\label{sec:sLSTM}

In order to enhance LSTMs by granting them the capacity to modify how information is stored, we incorporate exponential gates (depicted in red), along with mechanisms for normalization and stabilization. Specifically, the activation functions of input and forget gates are allowed to be exponential. To address normalization, we propose a normalizer state, which aggregates the result of the input gate multiplied by subsequent forget gates.

The scalar sLSTM forward pass is:

\begin{align}
\label{eq:slstmforward}
c_t \ &= \  f_t \ c_{t-1} \ + \ i_t \ z_t & & 
 &\text{cell state} \\
n_t \ &= \  f_t \ n_{t-1} \ + \ i_t  & & 
 &\text{normalizer state} \\
h_t \ &= \ o_t \ \tilde{h}_t \ ,
  &\tilde{h}_t \ &= \ c_t / n_t
&\text{hidden state} \\
z_t \ &= \ \varphi \left( \tilde{z}_t \right) \ , 
  &\tilde{z}_t \ &=  \ \mathbf{w}^\top_{z} \ \mathbf{x}_t \ + \
  r_{z}  h_{t-1} \ + \  b_{z}
  &\text{cell input} \\
i_t \ &= \ \!\exp \left( \tilde{i}_t  \right) \ , 
  &\tilde{i}_t \ &= \ \mathbf{w}^\top_{i} \ \mathbf{x}_t \ + \
  r_{i}  \ h_{t-1} \ + \  b_{i}
  &\text{input gate} \\
f_t \ &= \ \sigma \left(  \tilde{f}_t \right) \ \text{OR} \
\!\exp \left(  \tilde{f}_t \right) \ ,
  &\tilde{f}_t \ &= \ \mathbf{w}^\top_{f} \ \mathbf{x}_t  \ + \
  r_{f}  \ h_{t-1} \ + \  b_{f}
  &\text{forget gate} \\
o_t \ &= \ \sigma \left( \tilde{o}_t \right) \ , 
  &\tilde{o}_t  \ &= \ \mathbf{w}^\top_{o} \ \mathbf{x}_t \ + \
  r_{o}  \ h_{t-1} \ + \  b_{o}
  &\text{output gate} 
\end{align}

We transfer the original LSTM gating techniques, i.e., input- and/or hidden-dependent gating plus bias term, to the new architectures. Exponential activation functions can lead to large values that cause overflows. Therefore, we stabilize gates with an additional state \mbox{$m_t$~\citep{Milakov:18arxiv}}:

\begin{align}
\label{eq:slstmstabil}
m_t \ &= \ \max \left( \log ( f_t ) + m_{t-1} , \log ( i_t ) \right) &\text{stabilizer state} \\
i'_t \ &= \ \!\exp \left( \log \left ( i_t \right) - m_t \right) \ = \ \!\exp \left( \tilde{i}_t - m_t \right) \ \, 
  &\text{stabil. input gate} \\
f'_t \ &= \ \!\exp \left( \log \left( f_t \right) + m_{t-1} - m_t \right) \ \, 
  &\text{stabil. forget gate}
\end{align}

As demonstrated in Appendix~\ref{sec:appsLSTM}, substituting $f_t$ with $f'_t$ and $i_t$ with $i'_t$ during the network’s forward computation does not impact either the network's overall output or the gradients of the loss function relative to the model parameters.

\paragraph{New Memory Mixing.} The sLSTM architecture, similar to the conventional LSTM (refer to Appendix~\ref{sec:appsLSTM}), may incorporate several memory cells. Having multiple memory cells allows for memory mixing facilitated by recurrent connections $\mathbf{R}_{\mathbf{z}}$, $\mathbf{R}_{\mathbf{i}}$, $\mathbf{R}_{\mathbf{f}}$, $\mathbf{R}_{\mathbf{o}}$ 
which link the hidden state vector $\mathbf{h}$ to the inputs for the memory cell $\mathbf{z}$ and the gating mechanisms
$\mathbf{i}$, $\mathbf{f}$, $\mathbf{o}$, respectively. In this revised memory mixing scheme, exponential gating introduces a novel element. The updated sLSTM supports several heads, enabling memory mixing within individual heads, but not between different heads. The combination of head structure in sLSTM and exponential gating creates a distinctive approach to memory mixing.

\subsection{mLSTM}
\label{sec:mLSTM}

To boost the memory capacity of LSTMs, we substitute the traditional scalar memory cell $c \in \mathbb{R}$ with a matrix representation $\mathbf{C} \in \mathbb{R}^{d \times d}$. This allows retrieval operations to be executed through matrix multiplication. At each time step $t$, a pair of vectors is stored: the key $\mathbf{k}_t \in \mathbb{R}^d$ and the value $\mathbf{v}_t \in \mathbb{R}^d$, in line with terminology from the Transformer architecture. Subsequently, at time $t + \tau$, the stored value $\mathbf{v}_t$ should be accessed using a query vector $\mathbf{q}_{t+\tau} \in \mathbb{R}^d$. This approach corresponds to the framework of Bidirectional Associative Memories (BAMs) \citep{Kohonen:72,Anderson:72,Nakano:72,Anderson:77}.

The covariance update rule~\citep{Sejnowski:77,Dayan:91} for storing a key-value pair is

\begin{align}
\mathbf{C}_t \ &= \ \mathbf{C}_{t-1} \ + \ \mathbf{v}_{t} \ \mathbf{k}_t^\top \ .
\end{align}

We posit the application of layer normalization prior to the input projection into keys and values, resulting in vectors with zero mean. According to~\citep{Dayan:91}, the covariance update mechanism yields an optimal solution for enhancing the separability of binary vectors during retrieval—this is directly tied to achieving a maximal signal-to-noise ratio. Enhanced separability can be attained by confining the retrieval process to pairwise associations and accepting the attendant quadratic computational demands seen in attention mechanisms~\citep{Krotov:16,Krotov:17,Ramsauer:21}. Notably, the covariance update mechanism coincides with the Fast Weight Programmers model~\citep{Schmidhuber:92ncfastweights,Schlag:21}, which was subsequently modified to introduce a constant decay factor for scaling $\mathbf{C}_{t-1}$ and a fixed learning rate applied to the outer product term $\mathbf{v}_{t} \mathbf{k}_t ^\top$~\citep{Ba:16}. Following this conceptual framework, we embed the covariance-based update procedure within the architecture of an LSTM, drawing parallels between the forget gate and decay rate, the input gate and learning rate, and the scaling of the output vector by the output gate.

For this matrix memory, the normalizer state is the weighted sum of key vectors, where each key vector is weighted by the input gate and all future forget gates. Again, the normalizer state keeps record of the strength of the gates. Since the dot product between query and normalizer state can be close to zero, we use the absolute value of this dot product and lower bound it by a threshold (typically 1.0) as done previously~\citep{Sun:23arxiv}. The mLSTM forward pass is:

\begin{align}
\label{eq:mlstm_recurrent_begin}
\mathbf{C}_t \ &= \  f_t \ \mathbf{C}_{t-1} \ + \ 
  i_t \ \mathbf{v}_t \ \mathbf{k}_t^\top & &  &\text{cell state} \\
\mathbf{n}_t \ &= \  f_t \ \mathbf{n}_{t-1} \ + \ 
  i_t \ \mathbf{k}_t & &  &\text{normalizer state} \\
\mathbf{h}_t  \ &= \ \mathbf{o}_t \ \odot \ \tilde{\mathbf{h}}_t
\ , \qquad \qquad 
\tilde{\mathbf{h}}_t \ = \   \mathbf{C}_t \mathbf{q}_t \ / \ 
  \max \left\{ |\mathbf{n}_t^\top \mathbf{q}_t|, 1 \right\} 
   &&&\text{hidden state} \\
\mathbf{q}_t \ &= \ \mathbf{W}_q \ \mathbf{x}_t \ + \ \mathbf{b}_q  & &  &\text{query input} \\
\mathbf{k}_t \ &= \ \frac{1}{\sqrt{d}} \mathbf{W}_k \ \mathbf{x}_t \ + \ \mathbf{b}_k  & &  &\text{key input} \\
\mathbf{v}_t \ &= \ \mathbf{W}_v \ \mathbf{x}_t \ + \ \mathbf{b}_v  & &  &\text{value input} \\
i_t \ &= \ \!\exp \!  \! \left( \tilde{i}_t  \right) \ , 
 \qquad \qquad \ \ \ \ \,
  \tilde{i}_t \ = \ \mathbf{w}^\top_{i} \ \mathbf{x}_t \ + \  b_{i} \ \ \ 
  &&&\text{input gate} \\
f_t \ &= \ \sigma \!  \left(  \tilde{f}_t \right) \ \text{OR} \ \!\exp \!  \! \left(  \tilde{f}_t \right) , \ \ \: \!
\tilde{f}_t \ = \ \mathbf{w}^\top_{f} \ \mathbf{x}_t  \ + \
  b_{f}
  &&& \text{forget gate} \\
\label{eq:mlstm_recurrent_end}
\mathbf{o}_t \ &= \ \sigma \left( \tilde{\mathbf{o}}_t \right) \ , \qquad \qquad \quad \ \ \ \,
  \tilde{\mathbf{o}}_t  \ = \ \mathbf{W}_{\mathbf{o}} \ \mathbf{x}_t \ + \
  \mathbf{b}_{\mathbf{o}}
  &&&\text{output gate} 
\end{align}

Like the standard LSTM, mLSTM architecture supports several memory cells. In the context of mLSTM, utilizing multiple heads or multiple cells results in analogous behavior, due to the absence of memory mixing. To ensure stable operation of the exponential gates in mLSTM, we employ the stabilization approaches previously outlined for sLSTM (refer to Equation~\ref{eq:slstmstabil}). Owing to the lack of memory mixing in mLSTM, its recurrence relations lend themselves to a parallelized implementation. Additional information can be found in Appendix~\ref{sec:appmLSTM}.

\subsection{xLSTM Architecture}
\label{sec:xLSTMarch}

\paragraph{xLSTM Blocks.}
\label{sec:xLSTMblock}
The purpose of an xLSTM block is to achieve a nonlinear summary of previous states within a high-dimensional representation, thereby improving the discrimination between distinct histories or contexts. Accurate prediction of the next item in a sequence, such as the subsequent token, relies on the ability to separate such histories. Our approach utilizes Cover’s Theorem~\citep{Cover:65}, which asserts that nonlinear embeddings of patterns into higher-dimensional spaces enhance the likelihood of achieving linear separability compared to the original lower-dimensional space. We analyze two types of residual block designs: (i) A residual block featuring post up-projection (inspired by the Transformer architecture), which first performs nonlinear summarization in the original feature space, projects the result linearly into a higher-dimensional space, applies a nonlinear activation, and subsequently linearly projects back to the original feature space; see the leftmost illustration in Figure~\ref{fig:Backbones} and the third panel of Figure~\ref{fig:xlstm_sketch}. A comprehensive schematic can also be found in Figure~\ref{fig:appsLSTM_detailed} within the appendix. (ii) A residual block adopting pre up-projection (analogous to State Space Models), where a linear mapping into a high-dimensional space is applied at the outset,

compresses past information in a non-linear manner within a high-dimensional space and subsequently projects it back to the original dimensionality through a linear transformation. If an xLSTM block employs an sLSTM, the up-projection occurs after the sLSTM. Conversely, for xLSTM blocks utilizing an mLSTM, the up-projection precedes the mLSTM module, as the increase in dimensionality boosts memory capacity. Detailed illustrations can be found in the leftmost panel of Figure~\ref{fig:Backbones}, the third column in Figure~\ref{fig:xlstm_sketch}, or in Figure~\ref{fig:appsLSTM_detailed} within the appendix.

\paragraph{xLSTM Architecture. }
\label{sec:xLSTMarchitecure}
The xLSTM model is formed by stacking multiple blocks with residual connections~\citep{Srivastava:15,He:16}. This architecture adopts the prevalent pre-LayerNorm~\citep{Ba:16b} residual framework, consistent with those employed in modern Large Language Models. For visual reference, consult the final column of Figure~\ref{fig:xlstm_sketch}.

\subsection{Memory and Speed Considerations}

In contrast to Transformers, xLSTM architectures exhibit linear time complexity and maintain a fixed memory requirement independent of sequence length. Owing to their compressive memory design, xLSTM models are highly appropriate for practical, large-scale scenarios as well as for deployment on edge devices.

The memory component in mLSTM is parameter-free but imposes significant computational overhead due to its $d \times d$ matrix structure and the associated $d \times d$ updates. This represents a trade-off between available memory capacity and computational demands. However, since these operations can be parallelized across GPUs, they exert only a negligible influence on overall runtime.

Although mLSTM offers parallelization capabilities similar to FlashAttention~\citep{Dao:22,Dao:23} and GLA~\citep{Yang:23arxiv}, sLSTM does not support parallel execution because of its memory mixing via hidden-to-hidden interactions. To address this limitation, we have engineered an efficient CUDA-based implementation optimized down to the GPU register level, allowing sLSTM to run at speeds typically less than twice as slow as mLSTM.

\section{Related Work}
\label{sec:related}

\paragraph{Linear Attention.} To address the quadratic time complexity associated with context length in the Transformer architecture, various strategies have been devised to render attention mechanisms linear with respect to sequence length. The Synthesizer framework creates synthetic attention weights independently of token--token dependencies~\citep{Tay:20arxiv}. Linformer achieves self-attention through low-rank projection matrices, facilitating a linear approximation~\citep{Wang:20arxiv}. The Linear Transformer modifies the core attention mechanism so that it operates linearly~\citep{Katharopoulos:20}. Performer employs a positive orthogonal random features technique to produce a linearized approximation of the attention softmax~\citep{Choromanski:21}. Additionally, structured global convolutions, such as those in Structured Global Convolution (SGConv)~\citep{Li:22} and the Hyena Hierarchy~\citep{Poli:23}, have been proposed as substitutes for traditional attention, offering rapid computation for longer sequences.

\paragraph{State Space Models.} State Space Models (SSMs) have attracted significant attention lately due to their linear computational complexity with respect to sequence length, as well as their competitive results against Transformer-based methods. Early in this development, the Structured State Space sequence model (S4)~\citep{Gu:21} emerged, followed by advances such as the Diagonal State Space (DSS) model~\citep{Gupta:22}, Gated State Space (GSS) models~\citep{Mehta:22}, S5 model~\citep{Smith:22}, Bidirectional Gated SSM (BiGS)~\citep{Wang:22}, H3 model~\citep{Fu:23}, and Mamba~\citep{Gu:24arxiv}.

\paragraph{Recurrent Neural Networks.} In recent literature, Recurrent Neural Networks (RNNs) have been proposed as alternatives to Transformers and attention mechanisms, owing to their linear computational complexity with respect to sequence length. Language modeling tasks have benefited from RNNs equipped with Deep Linear Recurrent Units (LRUs), which have yielded encouraging outcomes~\citep{Orvieto:23, De:24arxiv}. Other variants, such as Hierarchically Gated Linear RNN (HGRN)~\citep{Qin:23} and its successor HGRN2~\citep{Qin:24arxiv}, have also demonstrated similar advancements. Among the RNN-based strategies for scaling language models, RWKV~\citep{Peng:23arxivshort,Peng:24arxivshort} stands out, delivering results that rival those achieved by Transformer architectures.

\paragraph{Gating.}
Gating stands out as a central concept introduced by LSTM, subsequently revisited and adapted across a variety of recent architectures. Numerous models, including HGRN~\citep{Qin:23}, HGRN2~\citep{Qin:24arxiv}, Gated Linear Attention (GLA)~\citep{Yang:23arxiv}, Gated State Space (GSS) models~\citep{Mehta:22}, Bidirectional Gated SSM (BiGS)~\citep{Wang:22}, Moving Average Equipped Gated Attention (MEGA)~\citep{Ma:22}, RWKV~\citep{Peng:23arxivshort}, and Mamba~\citep{Gu:24arxiv}, have implemented gating as a central mechanism within their respective frameworks.

\paragraph{Covariance Update Rule.} 
To bolster memory storage, the mLSTM unit was enhanced with a matrix-based memory governed by a covariance update rule. This principle is foundational in several other architectures employing similar update techniques, such as Fast Weight Programmers~\citep{Schmidhuber:92ncfastweights,Schlag:21}, RWKV-5 and RWKV-6~\citep{Peng:24arxivshort}, Retention~\citep{Sun:23arxiv}, Linear Transformer~\citep{Katharopoulos:20}, and HGRN2~\citep{Qin:24arxiv}.

\paragraph{Most Related.} xLSTM is most closely related, in terms of underlying concepts, to Retention~\citep{Sun:23arxiv}, RWKV~\citep{Peng:23arxivshort,Peng:24arxivshort}, and HGRN2~\citep{Qin:24arxiv}. These architectures incorporate matrix memory and gating mechanisms. Nevertheless, unlike the sLSTM, these methods lack the capability for memory mixing. This absence prevents them from effectively addressing state tracking challenges, a task that memory mixing facilitates. Consequently, LSTM architectures demonstrate a broader range of expressiveness compared to State Space Models (SSMs) and Transformers~\citep{Merrill:24,Deletang:23}, especially in applications such as code execution or monitoring entities within extended narratives, where accurate state tracking is essential.

\paragraph{Residually Stacking Architectures.} Modern large-scale deep learning models, including xLSTM, are formed by residual stacking of core modules~\citep{Srivastava:15,He:16}.

This design principle underpins the effectiveness of deep convolutional neural networks~\citep{He:16} as well as Transformers~\citep{Vaswani:17}. Transformers, in particular, constitute the foundational technology of today’s Large Language Models (LLMs) such as GPT-3~\citep{Brown:20short}, ChatGPT~\citep{Schulman:22short}, GPT-4~\citep{Achiam:23gpt4short}, Megatron-LM~\citep{Shoeybi:19}, Gopher~\citep{Rae:21short}, ERNIE 3.0 Titan~\citep{Wang:21arxivshort}, GLaM~\citep{Du:21glamshort}, Chinese M6~\citep{Lin:21}, AlexaTM 20B for multilingual tasks~\citep{Soltan:22}, OPT~\citep{Zhang:22opt}, Chinchilla~\citep{Hoffmann:22short}, BLOOM~\citep{Scao:22arxivshort}, GLM-130B~\citep{Zeng:22arxivshort}, LaMDA~\citep{Thoppilan:22short}, PaLM~\citep{Chowdhery:22short}, Llama~\citep{Touvron:23arxiv}, and Gemini~\citep{GeminiTeam:23arxiv,Reid:24arxiv}.

\section{Experiments}
\label{sec:Exp}

This section presents an empirical analysis of xLSTM, emphasizing its performance in language modeling relative to other approaches. Section~\ref{sec:ExpSyntethic} explores xLSTM’s effectiveness on a series of artificial benchmarks. In Section~\ref{sec:ExpComparison}, we evaluate the validation perplexity for a selection of modern language modeling techniques, all trained on a corpus of 15B tokens from SlimPajama~\citep{Soboleva:23}, and carry out ablation experiments focused on xLSTM. Subsequently, we investigate the scaling properties of each method, following methodologies similar to those described by \citet{Kaplan:20} and \citet{Brown:20short}. 

Section~\ref{sec:ExpLanguage} expands the scope with an in-depth language modeling study. Here, xLSTM and the top-performing models identified in Section~\ref{sec:ExpComparison} are further trained using 300B tokens from SlimPajama~\citep{Soboleva:23}. This evaluation proceeds in four parts: (1) measuring generalization to longer contexts, (2) analyzing both validation perplexity and task-specific performance on downstream tasks~\citep{Sutawika:24}, (3) assessing results across 571 domains included in the PALOMA language benchmark dataset~\citep{Magnusson:23arxivshort}, and (4) reexamining scaling properties, this time employing training sets that are twenty times larger than before.

\label{sec:xlstm_blocks_notation}
Throughout all experiments, we denote xLSTM[$a$:$b$] as representing an $a/b$ proportion between mLSTM-based and sLSTM-based xLSTM blocks. For instance, xLSTM[7:1] specifies that, from a total of eight blocks, seven utilize mLSTM while one employs sLSTM. When the total number of blocks is fixed at 48, this implies there are 6 sLSTM-based and 42 mLSTM-based blocks, respectively. Additionally, in every experimental setting, we include a pre up-projection block for mLSTM components and a post up-projection block for sLSTM components.

\subsection{Synthetic Tasks and Long Range Arena}
\label{sec:ExpSyntethic}
We begin by evaluating the impact of the novel exponential gating with memory mixing mechanism in xLSTM on formal languages~\citep{Deletang:23}. Subsequently, we investigate how the recently introduced matrix memory in xLSTM performs in the Multi-Query Associative Recall task~\citep{Arora:23arxiv}. Lastly, the ability of xLSTM to handle extensive sequences is analyzed within the framework of Long Range Arena~\citep{Tay:21}.

\paragraph{Evaluation of xLSTM's Exponential Gating and Memory Mixing Capabilities.} We evaluate xLSTM's novel exponential gating combined with memory mixing, which is hypothesized to facilitate effective state tracking, as suggested by prior work~\citep{Merrill:24, Merrill:23}. To do so, we build upon and enhance the formal language benchmarks introduced by \citet{Deletang:23}, expanding them to accommodate multi-length sequences in order to examine length extrapolation properties. Detailed methodology and comprehensive results for all tasks can be found in Appendix~\ref{sec:appExpSynthetic-formalLang}. For benchmarking purposes, we compare xLSTM against several architectures, including Transformers, State Space Models, and various Recurrent Neural Networks. Evaluation focuses on task-specific relevant tokens, with accuracy normalized between 0 (reflecting chance level) and 1 (indicating flawless performance). Specifically, we examine and contrast 2-block versions of the following models: xLSTM[0:1] (sLSTM only), xLSTM[1:0] (mLSTM only), xLSTM[1:1], Llama, Mamba, RWKV, Retention, Hyena, LSTM, and the variant integrating LSTM within Transformer blocks (LSTM (Block)). Results from these evaluations are presented in Figure~\ref{fig:formal-main}. Models such as Transformers and State Space Models, when used without memory mixing and thus lacking explicit state tracking capabilities, are unable to address tasks like recognizing regular grammars (e.g., the parity task). This limitation corroborates earlier findings that position both Transformers and State Space Models as inherently less capable in stateful sequence processing compared to RNNs~\citep{Merrill:24, Merrill:23, Deletang:23}.

\paragraph{Test of xLSTM's Memory Capacities on Associative Recall Tasks.} In this study, we assess the matrix-based memory of xLSTM by examining its memory capabilities on the Multi-Query Associative Recall task~\citep{Arora:23arxiv}. For each input sequence, key-value pairs are randomly drawn from an extensive vocabulary and must be stored for subsequent retrieval. To increase the task's complexity beyond its original design, we raise both the maximum number of key-value pairs to 256 and the sequence context length up to 2048 tokens, allowing for a more rigorous evaluation of model memory. We evaluate the performance of two-block versions of Llama, Mamba, RWKV-5, RWKV-6, xLSTM[1:1], and xLSTM[1:0]. Models are judged by their accuracy in recalling the correct pairs. Given that Transformer models like Llama possess memory scaling exponentially with the dimensionality of their representations~\citep{Ramsauer:21}, they serve as the benchmark for this challenge. Performance comparisons are illustrated in Figure~\ref{fig:mqar-main}. 
Among non-Transformer models, xLSTM[1:1] achieves the highest recall accuracy, even at lower model capacities. Notably, inclusion of the sLSTM block does not restrict memory ability; rather, it enhances it—particularly evident under the most demanding setting involving 256 key-value pairs. Supplementary findings appear in Appendix~\ref{sec:appExpSynthetic-mqar}, including extrapolation experiments that show xLSTM's improved memory mechanisms support generalization to longer contexts than those encountered in training.

\paragraph{Test of xLSTM's Long Context Capabilities on Long Range Arena.} In order to evaluate xLSTM’s effectiveness in managing extensive input sequences and broad contexts, we conduct a comparative study across various approaches using Long Range Arena~\citep{Tay:21}. xLSTM achieves robust results on each task, indicating that its architectural design facilitates notably efficient processing of diverse challenges associated with long context scenarios.

For more details, see Appendix~\ref{sec:appExpSynthetic-lra}.

\subsection{Method Comparison and Ablation Study}
\label{sec:ExpComparison}

In pursuit of the central inquiry of this study—namely, the capabilities of our novel LSTM variants in large-scale language modelling—we undertake training of xLSTMs, Transformers, State Space Models, among other approaches, utilizing 15B tokens from SlimPajama within an identical auto-regressive framework. The resulting models are evaluated on a shared validation set, and ablation analyses are conducted specifically for the xLSTM architectures.

\paragraph{Contrasting xLSTM with Alternative Approaches.} To provide a thorough comparison, we trained various models using 15B tokens sourced from SlimPajama~\citep{Soboleva:23}. These models were then assessed by measuring their validation perplexity. The methods included in our evaluation comprise: our novel xLSTM approach, GPT-3 (Transformer)~\citep{Brown:20short}, Llama (Transformer)~\citep{Touvron:23arxiv}, H3 (SSM)~\citep{Fu:23}, Mamba (SSM)~\citep{Gu:24arxiv}, RWKV-4, RWKV-5, and RWKV-6 (all RNNs)~\citep{Peng:23arxivshort, Peng:24arxivshort}, GLA (linear Transformer)~\citep{Yang:23arxiv}, HGRN and HGRN2 (RNNs)~\citep{Qin:23, Qin:24arxiv}, RetNet (linear Transformer)~\citep{Sun:23arxiv}, Hyena (linear Transformer)~\citep{Poli:23}, xLSTM[1:0], and xLSTM[7:1] (refer to Section~\ref{sec:xlstm_blocks_notation} for notation clarification). 
During training, mixed precision was employed; however, RWKV-5, RWKV-6, GLA, and HGRN2 did not utilize PyTorch's native automated mixed precision (see also Appendix Section~\ref{sec:appGenTrainProc}). We organize these approaches into three groups: (a) Transformers, (b) State Space Models (SSMs), and (c) Recurrent Neural Networks (RNNs), which include linear Transformers—these are variants that replace standard Transformer attention with linear mechanisms. Each model was designed to mirror the structure of a 350M parameter GPT-3, 
incorporating an embedding dimension of 1024 and 24 residual blocks. Notably, only GPT-3 leverages shared weights between token and output embeddings, resulting in a reduced parameter count relative to the others. As shown in Table~\ref{tab:spaj15b_model_results}, xLSTM demonstrates superior validation perplexity compared to all existing approaches. Further methodological specifics are available in Appendix~\ref{sec:appExpComparison}. 
Figure~\ref{fig:temp_sclaw15B} depicts the scaling dynamics observed in this experiment, suggesting that xLSTM's performance advantage persists as model size increases.

\paragraph{Ablation Studies.}
As presented in Table~\ref{tab:spaj15b_model_results} and Figure~\ref{fig:temp_sclaw15B}, the xLSTM model attains superior performance in language modeling when trained on 15B tokens sourced from SlimPajama. To systematically analyze the modifications transitioning from LSTM to xLSTM, we incrementally convert a conventional LSTM architecture into an xLSTM. The process begins with embedding LSTM layers within a pre-LayerNorm residual backbone. Next, this configuration is augmented by incorporating a post up-projection block. The final step entails the introduction of exponential gating along with matrix memory mechanisms.

The results are shown in Table~\ref{tab:ablstudies} (top). 

Through ablation experiments, we observe that the notable gains in performance are primarily due to the implementation of exponential gating as well as the inclusion of matrix memory. Recognizing the critical role of gating mechanisms in both RNNs and State Space Models, we further examine a variety of gating strategies. As shown in Table~\ref{tab:ablstudies} (bottom), results indicate that configuring each gate to be independently trainable and responsive to input leads to further improvements. Further detailed analyses regarding the specific backbone modules are provided in Appendix~\ref{sec:appAblDetails}.

\subsection{xLSTM as Large Language Model}
\label{sec:ExpLanguage}

In the final phase of our research, we conduct extensive language modeling trials to assess xLSTM’s viability as a large language model (LLM). To this end, we scale up the training dataset to encompass 300B tokens from SlimPajama, aligning with the token counts utilized in models such as Mamba~\citep{Gu:24arxiv} and Griffin~\citep{De:24arxiv}. 

Our evaluation includes a comparison between xLSTM and three baseline approaches: RWKV-4, Llama, and Mamba—each representing distinct categories of models referenced in Section~\ref{sec:ExpComparison}. RWKV-4 serves as the RNN comparator since reliable precision settings for RWKV-5, RWKV-6, and HGRN2 were only established after initiation of the 300B token runs (details provided in Appendix~\ref{sec:appGenTrainProc}). 

We explore a range of model capacities—specifically, 125M, 350M, 760M, and 1.3B parameters—and systematically examine their ability to generalize to longer sequences, their validation accuracy, and their effectiveness in various downstream applications. Their language modeling proficiency is measured across 571 text categories as defined in the PALOMA benchmark. Lastly, we analyze how these models adhere to scaling laws.

\paragraph{Sequence Length Extrapolation.}
\label{sec:spaj300B_ctx_extrapolation}
We begin by examining how well large-scale models (1.3B parameters) of xLSTM, RWKV-4, Llama, and Mamba extrapolate to longer sequences. Each model is trained using a context length of 2048 tokens, then evaluated on sequences with context lengths extended up to 16384. The outcomes are presented in Figure~\ref{fig:spaj300B_seq_extrapolation}. Notably, xLSTM sustains low perplexity scores at greater context lengths, outperforming the alternative approaches in this regard.

\paragraph{Validation Perplexity and Downstream Tasks.}
\label{sec:spaj_downstream_eval}
Next, across all parameter scales, we assess the performance of xLSTM, RWKV-4, Llama, and Mamba by measuring their perplexity on the SlimPajama validation split for next-token prediction, as well as their efficacy on downstream evaluations that assess capabilities in common sense reasoning.

The validation perplexities for various methods are reported in the third column of Table~\ref{tab:lmeval}. Across all model sizes, xLSTM[1:0] and xLSTM[7:1] consistently yield the lowest validation set perplexities, making them the optimal choices. The remaining columns in Table~\ref{tab:lmeval} summarize results on downstream tasks. xLSTM emerges as the leading method for most tasks and model configurations, with the exception of the ARC task where Mamba occasionally achieves superior performance. Additional information can be found in Appendix~\ref{sec:appExpLanguage}.

\paragraph{Performance on PALOMA Language Tasks.} As a third evaluation, we assess the next token prediction capabilities of xLSTM, RWKV-4, Llama, and Mamba across various model scales on PALOMA language tasks~\citep{Magnusson:23arxivshort}. The metric used is perplexity, computed for next token prediction across 571 different textual domains, encompassing a diverse set from nytimes.com to the r/depression subreddit.

Table~\ref{tab:ppleval} presents token prediction perplexity separated into two categories: language modeling tasks (displayed in the first seven columns) and detailed domain-specific benchmarks (featured in the last five columns). On these tasks, xLSTM[1:0] achieves superior results compared to xLSTM[7:1].

Specifically, xLSTM[1:0] yields lower perplexity than Mamba in 568 out of 571 domains (99.5\%), than Llama in 486 domains (85.1\%), and than RWKV-4 in 570 domains (99.8\%). Additional findings and explanations can be found in Appendix~\ref{sec:appExpLanguage}.

\paragraph{Scaling Laws.} In the fourth part of our analysis, we investigate power-law scaling dynamics, which facilitate the prediction of model performance as model size increases~\citep{Kaplan:20,Brown:20short}. The scaling trends are depicted in Figure~\ref{fig:temp_sclaw300B}. While the general scaling pattern is consistent across all model types, each exhibits a distinct offset. Among these, RWKV-4 yields the least favorable results, followed by Llama and Mamba, respectively. Notably, xLSTM achieves superior outcomes compared to Mamba, and the performance gap between xLSTM and Mamba is comparable to that between Mamba and Llama. The observed scaling behavior suggests that as models grow larger, xLSTM is likely to maintain its performance advantages relative to both Transformers and State-Space architectures.

\textbf{Generation Times and Maximal Throughput.} In our final set of evaluations, as illustrated in Figure~\ref{fig:inference_time_generation}, we assess the latency of text generation, alongside maximal throughput, presented in Figure~\ref{fig:inference_time_decoding_throughput} (left), focusing on the xLSTM architectures at the 1.3B parameter scale. For comparison, we include similarly sized implementations of Mamba, Llama, and RWKV sourced from HuggingFace; the Llama model utilizes a static key-value cache. It is notable that, during our investigation, neither the compilation for the entire cache of the Transformer Model nor \texttt{torch.compile}-based compilation for Mamba was operational. All models were evaluated in the text generation task with a batch size of 1 and a pre-fill value of 16 tokens, a setting particularly conducive to Transformer performance.

In Figure~\ref{fig:inference_time_generation}, results reveal that both xLSTM and the other recurrent baselines, Mamba and RWKV-4, demonstrate a linear relationship between generation time and input sequence length. In contrast, Llama displays quadratic growth in generation time. For the throughput experiment, we varied batch size and the number of prefilled tokens when running Llama. 

As seen in Figure~\ref{fig:inference_time_decoding_throughput} (right), the xLSTM maintains constant memory usage irrespective of batch size, thereby supporting substantially larger batches relative to Llama, and ultimately reaching the highest throughput among the evaluated models.

\section{Limitations}

(i) Unlike mLSTM, the approach to memory mixing in sLSTM restricts the possibility for parallel computation, rendering efficient parallel processing infeasible. Despite this limitation, we engineered a CUDA kernel for sLSTM which demonstrates reasonably competitive speed, running at less than twice the time of our parallel mLSTM implementation. (ii) The current CUDA kernels for mLSTM lack optimization, resulting in approximately four times slower performance compared to FlashAttention or the scan utilized in Mamba. There remains substantial potential to improve the speed of mLSTM by employing optimization strategies similar to those used in FlashAttention. (iii) The matrix memory of mLSTM requires processing $d \times d$ matrices, which increases computational complexity. However, as parameter-free memory update and retrieval can leverage typical matrix operations with parallel processing, the actual wall time increase from handling complex memory remains limited. (iv) Careful selection of initial values for the forget gates is necessary to ensure proper functionality. (v) The matrix memory of mLSTM does not depend on sequence length, but expanding the sequence may exhaust memory resources for significantly extended contexts. For contexts up to 16k tokens, this has not been a concern; refer to Section~\ref{sec:spaj300B_ctx_extrapolation} for empirical validation. (vi) Because large-scale language experiments are computationally demanding, neither the architectural designs nor hyperparameters—particularly for larger xLSTM variants—have been thoroughly optimized. We expect that unlocking the full capacity of xLSTM will require substantial and deliberate optimization efforts.

\section{Conclusion}
We have provided a partial answer to the question: To what extent does scaling LSTM architectures to billions of parameters advance language modeling? The evidence so far suggests: ``At least to the level achieved by contemporary models such as Transformers or State Space Models.'' By introducing exponential gating with memory mixing and redesigning memory structures, we developed the xLSTM variant. Empirical results demonstrate that xLSTM models are competitive with, and sometimes outperform, leading approaches such as Transformers and State Space Models in language modeling tasks. Scaling analysis further implies that larger xLSTM configurations could represent formidable alternatives to existing Transformer-based Large Language Models. Furthermore, xLSTM holds promise for applications beyond language modeling, potentially influencing areas including Reinforcement Learning, Time Series Prediction, and physical system modeling.

\end{document}